\par\noindent\textbf{(1)} 单粒子 \(z=1+e^{-\beta\Delta}\)，总 \(Z=z^N\)。激发态概率 \(p=e^{-\beta\Delta}/z=1/(1+e^{\beta\Delta})\)，基态概率 \(1-p\)，故平均能量 \(U=N\Delta p\)。这是独立伯努利占有的乘积分布而非Bose无限占有。
\par\noindent\textbf{(2)} 令 \(x=\Delta/(k_BT)\)，热容 \(C_V=\partial_TU=Nk_Bx^2e^x/(1+e^x)^2=Nk_Bx^2/[4\cosh^2(x/2)]\)。能量方差 \(\operatorname{Var}(H)=N\Delta^2p(1-p)\)，可直接验证 \(\operatorname{Var}(H)=k_BT^2C_V\)。
\par\noindent\textbf{(3)} 对无量纲热容 \(c(x)=x^2/[4\cosh^2(x/2)]\) 求导，得到 \(c'/c=2/x-\tanh(x/2)\)。正极值满足 \(x\tanh(x/2)=2\)，唯一正根 \(x_*\approx2.39936\)，因此峰温 \(T_*=\Delta/(k_Bx_*)\)，峰值有限而不代表相变。
\par\noindent\textbf{(4)} 低温 \(x\gg1\) 时 \(C_V\sim Nk_Bx^2e^{-x}\) 指数抑制；高温 \(x\ll1\) 时 \(C_V\sim Nk_Bx^2/4\to0\)。若整体平移两个能级，\(U\) 增常量但热容、概率及方差不变，说明零点选择不影响可测热容。
